% !TeX root = ../../tesis.tex
% Esquema ilustrativo (decisión del autor 2026-11-04): claves y rutas
% tomadas de VFTModel/src/api/main.py; valores = marcadores de tipo.

\section{Ejemplo de Respuesta: Fuerza Capilar \texorpdfstring{($FC$)}{(FC)}}
\label{anx:vft-out-capilar}

Esquema ilustrativo de la respuesta del indicador $FC$
(§\ref{subsec:alimentacion-capilar}). Los nombres de las claves y la ruta son los
de \texttt{VFTModel}; los valores son marcadores de tipo. Las cifras de la
corrida 2026-10-04 se reportan en el Capítulo~\ref{ch:analisis-red-actual}.

Los registros vienen ordenados de mayor a menor \texttt{Fuerza\_Capilar\_Total}.

\begin{lstlisting}[style=estiloPython,
  caption={Esquema de respuesta del endpoint
           \texttt{/topological/capillary-strength} (ilustrativo).},
  label=lst:out-capilar]
{
  "status": "success",
  "parametros": { "modo_grafo": "REALISTIC_INTEGRATION",
                  "tolerancia_transbordo_grafo_m": <número>,
                  "tolerancia_agrupacion_algoritmo_m": <número> },
  "data": [
    { "Nodo_ID": "(<lon>, <lat>)",
      "Estacion": "<estación>",
      "Sistemas_Integrados": ["<sistema>"],
      "Detalle_Estaciones": { "<sistema>": ["<estación>"] },
      "Estaciones_Agrupadas": 1,
      "Conexiones_Entrada": <número>,
      "Conexiones_Salida": <número>,
      "Fuerza_Capilar_Total": <número> }
  ]
}
\end{lstlisting}
